\begin{center}\begin{tikzpicture}[font=\small,>=stealth,line join=round]
\draw[gray!40,rounded corners=3pt] (0,0) rectangle (14.2,8.3);\node at (7.1,7.700000000000001) {奇函数绕原点对称：左、中、右连成同向单调};
\begin{scope}[shift={(4.1,3.9499999999999997)},x=1.3636363636363635cm,y=0.32857142857142857cm]\draw[->,gray] (-2.2,0)--(2.2,0) node[right]{$x$};\draw[->,gray] (0,-7)--(0,7) node[above]{$y$};\draw[blue,thick,solid,domain=-1.8:1.8,samples=180] plot(\x,{(\x)^3});\fill[orange!85!black] (-1,-1) circle(2.2pt);\fill[orange!85!black] (0,0) circle(2.2pt);\fill[orange!85!black] (1,1) circle(2.2pt);\end{scope}\node[anchor=west,align=left,text width=5.0cm] at (8.4,4.0) {图取 $f(x)=x^3$\\[9pt]左侧 $f(u)<0$\\[9pt]原点 $f(0)=0$\\[9pt]右侧 $f(v)>0$\\[9pt]跨越 $0$ 仍有 $f(u)<f(v)$};
\node[text width=12.6cm,align=center] at (7.1,.65) {证明还要比较同在左侧的两点；用奇性将它们对应到右侧即可。};
\end{tikzpicture}\end{center}
